\documentclass[aspectratio=169]{beamer}
% Préambule Beamer commun aux huit séquences.
% Les decks restent volontairement compatibles avec Beamer standard.

\usepackage[T1]{fontenc}
\usepackage[utf8]{inputenc}
\usepackage[french]{babel}
\usepackage{graphicx}
\usepackage{booktabs}
\usepackage{hyperref}

\setbeamertemplate{navigation symbols}{}
\setbeamertemplate{footline}[frame number]

\newcommand{\decisionquestion}[1]{%
  \begin{block}{Question décisionnelle}
    #1
  \end{block}
}

\newcommand{\confidencelevel}[1]{\textbf{Niveau de confiance :} #1}

\graphicspath{{figures/}}
\title{De la donnée capteur à la décision opérationnelle}
\subtitle{Séquence 7 -- Confiance, sécurité et robustesse de la chaîne}
\author{École de l'air et de l'espace}
\date{}

\begin{document}

\begin{frame}
  \titlepage
\end{frame}

\begin{frame}{Bienvenue dans la séquence 7}
  \begin{block}{L'idée qui traverse toute la séance}
    Les séquences précédentes ont supposé que chaque message provenait bien du capteur annoncé, envoyé une seule fois, au moment annoncé. Cette séquence retire cette hypothèse et demande : que vaut un indicateur parfaitement calculé si la chaîne qui l'alimente ne peut pas être qualifiée de confiance ?
  \end{block}
  \begin{itemize}
    \item Le matin, nous distinguons ce qu'une authentification et une liste de contrôle d'accès (ACL) protègent réellement, et ce qu'elles ne protègent pas.
    \item L'après-midi, nous analysons un lot dont certains messages sont exactement ce qu'ils prétendent être -- et d'autres non -- sans qu'aucune valeur ne soit absurde en elle-même.
    \item La question n'est jamais « ce message est-il vrai ? », mais « quelle confiance cette chaîne mérite-t-elle, et quelle décision cette confiance permet-elle encore ? »
  \end{itemize}
\end{frame}

\begin{frame}{Le fil rouge des huit séquences}
  \centering
  \Large Donnée capteur imparfaite\\[2mm]
  $\Downarrow$\\[-1mm]
  Pipeline de traitement\\[-1mm]
  $\Downarrow$\\[-1mm]
  Indicateur, visualisation\\[-1mm]
  $\Downarrow$\\[-1mm]
  \alert{Confiance dans la chaîne elle-même}\\[-1mm]
  $\Downarrow$\\[-1mm]
  Décision opérationnelle
  \vfill
  \normalsize Les séquences 1 à 6 ont construit, nettoyé, résumé et présenté la donnée. Cette séquence questionne une étape antérieure à toutes les autres : la chaîne qui transporte cette donnée mérite-t-elle d'être crue ?
\end{frame}

\begin{frame}{Ce que vous saurez faire à la fin de cette séquence}
  \begin{enumerate}
    \item définir intégrité, authentification, ACL, injection et rejeu, et relier chacun à un effet décisionnel précis ;
    \item comparer un broker ouvert et un broker protégé sur des cas concrets d'accès autorisé et refusé ;
    \item identifier les vulnérabilités d'une chaîne MQTT à partir d'une configuration donnée ;
    \item analyser un lot suspect pour y détecter doublons, rejeux et ruptures, sans aucun nouveau calcul ;
    \item classer des hypothèses concurrentes par probabilité et impact, et recommander si l'on peut agir sur ces données.
  \end{enumerate}
\end{frame}

\begin{frame}{La question qui guidera la séance}
  \decisionquestion{Que vaut une décision opérationnelle si l'on ne peut pas qualifier la confiance dans la chaîne de données ?}
  \begin{itemize}
    \item Il ne s'agit pas de prouver qu'une chaîne est parfaitement sûre -- aucune ne l'est -- mais de savoir nommer précisément ce qu'elle protège et ce qu'elle ne protège pas.
    \item Un message peut être numériquement plausible et pourtant ne pas mériter d'être suivi.
    \item Le raisonnement attendu distingue la confiance dans la source de la confiance dans le traitement qui la reçoit.
  \end{itemize}
\end{frame}

\begin{frame}[plain]
  \centering
  \Huge Section 1\\[3mm]
  \Large Un flux qui ne se comporte pas normalement
  \vfill
  \normalsize Avant tout concept, un simple constat : quelque chose dans ce lot ne va pas.
\end{frame}

\begin{frame}{Le scénario : des messages incohérents}
  \begin{columns}[T]
    \column{.55\textwidth}
    Un extrait de `comms-shelter-01` contient un message dont l'horodatage de mesure est postérieur à sa réception, et un autre qui semble revenir sur une mesure déjà vue vingt-neuf minutes plus tôt.
    \column{.4\textwidth}
    \begin{block}{Contrainte de décision}
      Aucune de ces deux anomalies ne porte sur la valeur physique elle-même : la température reste dans une plage parfaitement normale.
    \end{block}
  \end{columns}
  \vfill
  \alert{La première question n'est donc pas « la mesure est-elle plausible ? » mais « ce message est-il bien ce qu'il prétend être ? »}
\end{frame}

\begin{frame}{Avant tout outil : votre décision}
  \begin{columns}[T]
    \column{.5\textwidth}
    Sur la seule base de ce constat, choisissez :
    \begin{enumerate}
      \item[A.] traiter le lot normalement, rien ne semble alarmant en valeur ;
      \item[B.] isoler le lot avant toute analyse plus poussée ;
      \item[C.] traiter les zones non concernées, isoler `comms-shelter-01` ;
      \item[D.] demander une vérification terrain avant tout traitement.
    \end{enumerate}
    \column{.42\textwidth}
    Notez également :
    \begin{itemize}
      \item votre confiance de 0 à 100 \% ;
      \item ce qui, dans un message MQTT, pourrait ne pas être fiable ;
      \item ce que la seule valeur mesurée ne peut jamais garantir.
    \end{itemize}
  \end{columns}
\end{frame}

\begin{frame}[plain]
  \centering
  \Huge Section 2\\[3mm]
  \Large Intégrité, authentification, ACL, injection, rejeu
  \vfill
  \normalsize Cinq notions, puis une démonstration : le même broker, avec et sans protection.
\end{frame}

\begin{frame}{Intégrité}
  \begin{block}{Définition}
    La garantie qu'un message reçu est identique à celui qui a été émis, sans altération en chemin.
  \end{block}
  \begin{itemize}
    \item \textbf{Piège fréquent :} confondre intégrité et exactitude -- un message intègre peut transporter une mesure fausse dès l'origine.
    \item \textbf{Effet sur la décision :} sans garantie d'intégrité, un message correctement reçu n'est pas nécessairement le message réellement envoyé par le capteur.
  \end{itemize}
\end{frame}

\begin{frame}{Authentification}
  \begin{block}{Définition}
    La vérification de l'identité d'un client avant de le laisser se connecter au broker.
  \end{block}
  \begin{itemize}
    \item \textbf{Piège fréquent :} croire qu'un broker accessible uniquement en local (\texttt{127.0.0.1}) n'a pas besoin d'authentification.
    \item \textbf{Effet sur la décision :} sans authentification, n'importe quel client capable d'atteindre le port peut publier un message qui semble provenir d'un capteur légitime.
  \end{itemize}
\end{frame}

\begin{frame}{Liste de contrôle d'accès (ACL)}
  \begin{block}{Définition}
    Un ensemble de permissions précises, par utilisateur et par topic, appliquées après l'authentification.
  \end{block}
  \begin{itemize}
    \item \textbf{Piège fréquent :} croire que l'authentification seule suffit -- un utilisateur authentifié mais non restreint peut publier ou lire n'importe quel topic.
    \item \textbf{Effet sur la décision :} une ACL bien conçue limite un compte capteur compromis à son seul périmètre ; son absence transforme tout compte en accès total.
  \end{itemize}
\end{frame}

\begin{frame}{Injection}
  \begin{block}{Définition}
    L'introduction dans le flux d'un message qui n'a jamais été produit par un capteur légitime.
  \end{block}
  \begin{itemize}
    \item \textbf{Piège fréquent :} supposer qu'un message injecté sera forcément une valeur absurde, donc facile à repérer.
    \item \textbf{Effet sur la décision :} un message injecté avec une valeur plausible peut influencer une décision exactement comme un message authentique.
  \end{itemize}
\end{frame}

\begin{frame}{Rejeu (replay)}
  \begin{block}{Définition}
    La republication d'un message déjà vu, sous une identité de message différente, pour le faire passer pour une nouvelle mesure.
  \end{block}
  \begin{itemize}
    \item \textbf{Piège fréquent :} chercher un rejeu uniquement parmi les messages strictement identiques -- un rejeu change délibérément l'identifiant ou le topic.
    \item \textbf{Effet sur la décision :} un rejeu peut faire réapparaître une mesure ancienne, rassurante ou alarmante, à un moment où elle ne reflète plus la situation réelle.
  \end{itemize}
  \decisionquestion{Un contrôle qui vérifie qu'un horodatage n'est pas dans le futur suffit-il à détecter un rejeu ?}
\end{frame}

\begin{frame}{Auto-vérification avant la démonstration}
  \small
  \begin{block}{Un broker accessible uniquement en local peut-il quand même être attaqué ?}
    Oui : toute machine ou tout conteneur qui atteint ce port -- y compris un processus compromis sur la même machine -- peut s'y connecter si aucune authentification n'est exigée.
  \end{block}
  \begin{block}{L'authentification suffit-elle à empêcher un abus de droits ?}
    Non : un compte authentifié mais sans ACL peut publier ou lire sur n'importe quel topic, y compris ceux d'autres capteurs ou d'autres zones.
  \end{block}
\end{frame}

\begin{frame}[fragile]{Démonstration -- broker ouvert vs broker protégé}
  \scriptsize
  \begin{verbatim}
docker compose -f docker/docker-compose.yml up -d --wait mosquitto mosquitto-protected

# 1. anonyme sur le broker ouvert (1883) : reussit
mosquitto_pub -h 127.0.0.1 -p 1883 -t airbase/test -m hello

# 2. anonyme sur le broker protege (1884) : refuse a la connexion
mosquitto_pub -h 127.0.0.1 -p 1884 -t airbase/test -m hello
  \end{verbatim}
  \normalsize
  Avant d'exécuter, prédisez : la seconde commande échoue-t-elle silencieusement, ou avec un message explicite ?
\end{frame}

\begin{frame}{Ce que révèle la démonstration}
  \tiny
  \begin{center}
  \begin{tabular}{@{}lll@{}}
    \toprule
    Tentative & Résultat & Visible pour l'émetteur ? \\
    \midrule
    anonyme $\rightarrow$ ouvert (1883) & acceptée & -- \\
    anonyme $\rightarrow$ protégé (1884) & refusée à la connexion & \textbf{oui, explicitement} \\
    \texttt{capteur-lora} publie & acceptée (ACL autorisée) & -- \\
    \texttt{superviseur} publie & \textbf{rejetée par l'ACL} & \textbf{non -- silencieux} \\
    \texttt{capteur-lora} s'abonne & connecté, \textbf{rien reçu} & \textbf{non -- silencieux} \\
    \bottomrule
  \end{tabular}
  \end{center}
  \vfill
  \alert{Un mauvais mot de passe ou une connexion anonyme sur un broker protégé échoue bruyamment. Un abus de droits sur un compte pourtant authentifié échoue silencieusement : la commande de \texttt{superviseur} se termine sans erreur, alors que son message n'a jamais été délivré.}
\end{frame}

\begin{frame}{Exercice -- identifier les vulnérabilités d'une chaîne MQTT}
  \begin{block}{Votre mission}
    À partir de la configuration observée, listez les vulnérabilités présentes et celles que l'ACL protégée corrige.
  \end{block}
  \begin{enumerate}
    \item le broker ouvert de \texttt{mosquitto.conf} accepte-t-il des clients anonymes ? Qu'est-ce que cela permet à un tiers non autorisé ?
    \item une ACL empêche-t-elle un compte compromis de lire les topics d'une autre zone ?
    \item les mots de passe et les messages transitent-ils chiffrés entre le client et le broker dans cette configuration ?
    \item un message parfaitement authentifié et autorisé peut-il quand même être un rejeu ?
  \end{enumerate}
  \begin{block}{Livrable}
    Une liste de vulnérabilités, chacune associée à la protection qui la couvre, la protège partiellement, ou ne la couvre pas du tout dans ce laboratoire.
  \end{block}
\end{frame}

\begin{frame}{Débrief -- ce que l'ACL protège, ce qu'elle ne protège pas}
  \begin{columns}[T]
    \column{.47\textwidth}\textbf{Couvert par l'authentification et l'ACL}
    \begin{itemize}
      \item connexion anonyme non autorisée ;
      \item lecture ou écriture hors du périmètre d'un compte.
    \end{itemize}
    \column{.47\textwidth}\textbf{Non couvert dans ce laboratoire}
    \begin{itemize}
      \item absence de chiffrement (TLS) entre client et broker ;
      \item rejeu d'un message par un compte pourtant autorisé à publier ;
      \item compromission du compte lui-même (mot de passe volé).
    \end{itemize}
  \end{columns}
  \vfill
  \decisionquestion{Une chaîne protégée par authentification et ACL est-elle une chaîne digne de confiance ?}
\end{frame}

\begin{frame}[plain]
  \centering
  \Huge Section 3\\[3mm]
  \Large Confiance dans la source, confiance dans le traitement
  \vfill
  \normalsize Deux confiances distinctes, à qualifier séparément avant toute décision.
\end{frame}

\begin{frame}{Confiance dans la source}
  \begin{block}{Définition}
    Le degré de certitude que ce message provient bien du capteur, de la zone et de l'instant qu'il annonce.
  \end{block}
  \begin{itemize}
    \item \textbf{Piège fréquent :} l'assimiler à la plausibilité de la valeur mesurée -- une source usurpée peut transmettre une valeur parfaitement crédible.
    \item \textbf{Effet sur la décision :} une source non qualifiée fragilise toute décision, même appuyée sur des indicateurs impeccablement calculés.
  \end{itemize}
\end{frame}

\begin{frame}{Confiance dans le traitement}
  \begin{block}{Définition}
    Le degré de certitude que la pipeline qui reçoit un message ne l'altère pas et ne masque pas ses anomalies.
  \end{block}
  \begin{itemize}
    \item \textbf{Piège fréquent :} croire qu'un contrôle qualité (séquence 4) couvre automatiquement les anomalies de sécurité -- un rejeu passe les contrôles de champs, d'unité et de plage sans aucune difficulté.
    \item \textbf{Effet sur la décision :} un traitement qui ne croise jamais l'identité des messages entre eux ne peut détecter un rejeu, quelle que soit la qualité de ses contrôles ligne à ligne.
  \end{itemize}
  \decisionquestion{Un message peut-il franchir tous les contrôles de qualité de la séquence 4 et rester un rejeu ?}
\end{frame}

\begin{frame}[fragile]{Démonstration -- diagnostiquer un lot suspect}
  \scriptsize
  \begin{verbatim}
python3 -m iot_decision.chain_trust_cli data/samples/batch004_suspect_scenario.jsonl
  \end{verbatim}
  \normalsize
  Ce module ne recalcule rien : il relit les signaux déjà détectables par les modules des séquences 3 et 4 (candidat de rejeu, doublon exact, incohérence temporelle, silence), et les relie à des hypothèses concurrentes.
\end{frame}

\begin{frame}{Activité -- Analyse d'un lot suspect}
  \begin{block}{Votre mission}
    Exécuter le diagnostic sur \texttt{batch004\_suspect\_scenario.jsonl} et interpréter chaque signal.
  \end{block}
  \begin{enumerate}
    \item relever le nombre de doublons exacts, de candidats de rejeu, d'incohérences temporelles et la durée du silence non expliqué ;
    \item pour chaque signal, identifier la ou les lignes brutes concernées dans le fichier JSONL ;
    \item vérifier que la mesure du candidat de rejeu passait bien tous les contrôles de qualité de la séquence 4 ;
    \item noter si un seul de ces signaux, pris isolément, aurait suffi à alerter.
  \end{enumerate}
\end{frame}

\begin{frame}{Ce que révèle le diagnostic}
  \small
  \begin{center}
  \begin{tabular}{@{}lc@{}}
    \toprule
    Signal & Valeur observée \\
    \midrule
    Doublons exacts & 1 \\
    Candidats de rejeu (même mesure, identité différente) & 1 \\
    Incohérences temporelles & 1 \\
    Silence non expliqué (maximum) & 15 min \\
    \bottomrule
  \end{tabular}
  \end{center}
  \vfill
  \normalsize Le candidat de rejeu réutilise la mesure de \texttt{comms-shelter-01} de 09:40 (29,3\,°C), republiée sous un nouvel identifiant trente minutes plus tard. Cette mesure est parfaitement valide au sens de la séquence 4 : aucun champ manquant, aucune unité incohérente, aucun horodatage postérieur à la réception.
\end{frame}

\begin{frame}[plain]
  \centering
  \Huge Section 4\\[3mm]
  \Large Décider malgré l'incertitude sur la chaîne
  \vfill
  \normalsize Quatre hypothèses concurrentes, une seule question : peut-on agir sur ces données ?
\end{frame}

\begin{frame}{Matrice des risques data/cyber}
  \small
  \begin{center}
  \begin{tabular}{@{}llll@{}}
    \toprule
    Hypothèse & Probabilité & Impact & Justification \\
    \midrule
    \textbf{Suspicion data/cyber} & \textbf{forte} & \textbf{élevé} & candidat de rejeu + incohérence temporelle \\
    Problème réseau & forte & moyen & silence non expliqué de 15 min \\
    Panne capteur & moyenne & moyen & 1 retransmission exacte détectée \\
    Incident réel & faible & élevé & aucun franchissement de seuil observé \\
    \bottomrule
  \end{tabular}
  \end{center}
  \vfill
  \normalsize Plusieurs hypothèses restent probables simultanément : elles ne s'excluent pas, elles orientent des vérifications différentes.
\end{frame}

\begin{frame}{Décision -- peut-on agir sur ces données ?}
  \decisionquestion{Que vaut une décision opérationnelle si l'on ne peut pas qualifier la confiance dans la chaîne de données ?}
  \begin{block}{Recommandation sur ce lot}
    Ne pas agir directement sur ces données : isoler le lot, vérifier l'identité et la source des messages suspects avant toute décision opérationnelle.
  \end{block}
  \normalsize Cette recommandation ne rejette pas le lot : elle en retarde l'usage jusqu'à vérification, exactement comme une confiance faible retardait une conclusion dans les séquences précédentes.
\end{frame}

\begin{frame}{Recommandations de sécurisation minimale}
  \begin{enumerate}
    \item désactiver l'accès anonyme dès qu'un broker dépasse le cadre d'un laboratoire strictement local ;
    \item attribuer une ACL par rôle, jamais un compte unique partagé entre tous les capteurs ;
    \item conserver l'identifiant de message et le topic d'origine pour pouvoir croiser les messages entre eux, pas seulement les valider un par un ;
    \item ne jamais déduire qu'un message est sûr du seul fait qu'il a franchi les contrôles de qualité de la séquence 4.
  \end{enumerate}
  \decisionquestion{Laquelle de ces quatre recommandations aurait, seule, empêché le candidat de rejeu de ce lot ?}
\end{frame}

\begin{frame}{Restitution -- 150 mots maximum}
  Rédigez une recommandation contenant :
  \begin{enumerate}
    \item la ou les hypothèses retenues, avec probabilité et impact ;
    \item un niveau de confiance dans la chaîne, distinct de la confiance dans les valeurs mesurées ;
    \item deux preuves chiffrées et retrouvables, avec leur fichier source ;
    \item une vérification prioritaire avant toute décision opérationnelle.
  \end{enumerate}
\end{frame}

\begin{frame}{Transition vers la suite du cours}
  Le mini-projet final s'appuiera directement sur cette séquence :
  \begin{enumerate}
    \item un lot ambigu peut combiner un problème de qualité (séquence 4), un indicateur trompeur (séquence 5) et un signal de confiance dans la chaîne (séquence 7) ;
    \item aucune de ces trois familles de vérification n'en remplace une autre.
  \end{enumerate}
  \vfill
  Retenez ce réflexe : un message qui franchit tous les contrôles de qualité n'a pas nécessairement franchi un contrôle de confiance dans la chaîne -- ce sont deux vérifications différentes.
\end{frame}

\begin{frame}{Activité -- Synthèse individuelle}
  Complétez les trois phrases, seul et sans écran :
  \begin{enumerate}
    \item Un message peut être valide au sens de la séquence 4 et pourtant ...
    \item L'authentification protège contre ... ; l'ACL protège en plus contre ...
    \item Avant d'agir sur un lot dont la chaîne n'est pas qualifiée, je vérifierais ...
  \end{enumerate}
  \vfill
  \alert{Message incohérent $\rightarrow$ signaux de confiance $\rightarrow$ hypothèses concurrentes $\rightarrow$ isolement et vérification $\rightarrow$ décision différée, pas refusée.}
\end{frame}

\begin{frame}{Ce que la séance a permis d'établir -- synthèse de référence}
  \small
  \begin{block}{Résultat de référence sur ce lot}
    Un doublon exact, un candidat de rejeu et une incohérence temporelle sur \texttt{comms-shelter-01}, plus un silence non expliqué de 15 minutes, placent « suspicion data/cyber » en tête des hypothèses (probabilité forte, impact élevé), devant un problème réseau plausible mais secondaire.
  \end{block}
  \begin{block}{Une recommandation défendable, à titre d'exemple}
    Isoler le lot, vérifier l'identité et la source des messages suspects avant toute décision opérationnelle ; ne jamais confondre la validité ligne à ligne d'un message avec la confiance dans son identité.
  \end{block}
  \normalsize
  \alert{D'autres formulations restent recevables si elles distinguent explicitement confiance dans la source, confiance dans le traitement, et qualité au sens de la séquence 4.}
\end{frame}

\begin{frame}{Références et ressources}
  \tiny
  \begin{itemize}
    \item Eclipse Mosquitto, documentation \texttt{mosquitto.conf} (authentification, \texttt{password\_file}, \texttt{acl\_file}) : \url{https://mosquitto.org/man/mosquitto-conf-5.html}
    \item OASIS, \emph{MQTT Version 5.0 Specification}, section sur les messages retained et la persistance : \url{https://docs.oasis-open.org/mqtt/mqtt/v5.0/mqtt-v5.0.html}
    \item NIST IR 8259A, \emph{IoT Device Cybersecurity Capability Core Baseline}, 2020 : \url{https://doi.org/10.6028/NIST.IR.8259A}
    \item OWASP, \emph{OWASP Internet of Things Top 10} : \url{https://owasp.org/www-project-internet-of-things/}
    \item ANSSI, recommandations de sécurité relatives aux objets connectés : \url{https://www.ssi.gouv.fr/}
  \end{itemize}
\end{frame}

\end{document}
